\documentclass[11pt,a4paper,modern]{FFExercises}
\usepackage[english,greek]{babel}
\usepackage[utf8]{inputenc}
\usepackage{nimbusserif}
\usepackage[T1]{fontenc}
\usepackage{amsmath}
\let\myBbbk\Bbbk
\let\Bbbk\relax
\usepackage[amsbb,subscriptcorrection,zswash,mtpcal,mtphrb,mtpfrak]{mtpro2}
\usepackage{graphicx,multicol,multirow,enumitem,tabularx,mathimatika,gensymb,venndiagram,hhline,longtable,tkz-euclide,fontawesome5,eurosym,tcolorbox,tabularray,tikzpagenodes,relsize}
\definecolor{xrwma}{HTML}{aa1212}
\usetikzlibrary{calc}
\usetikzlibrary{positioning}
\tcbuselibrary{skins,theorems,breakable}
\renewcommand{\textstigma}{\textsigma\texttau}
\renewcommand{\textdexiakeraia}{}

\ekthetesdeiktes
\begin{document}

\titlos{Μαθηματικά προσανατολισμού}{Β΄ Λυκείου}{Εσωτερικό γινόμενο}
\paragraph{Εσωτερικό γινόμενο (ορισμός)}

\askhsh
Δίνονται τα διανύσματα $\vec{a}, \vec{b}$ με $|\vec{a}|=4$, $|\vec{b}|=2$ και γωνία $(\widehat{\vec{a},\vec{b}})=\frac{\pi}{3}$. Να υπολογίσετε:
\begin{enumerate}
    \item Το εσωτερικό γινόμενο $\vec{a}\cdot\vec{b}$.
    \item Το εσωτερικό γινόμενο $\vec{b}\cdot\vec{a}$.
    \item Το τετράγωνο $\vec{a}^2$.
    \item Το γινόμενο $\vec{a}\cdot(2\vec{b})$.
\end{enumerate}

\askhsh
Σε τρίγωνο $AB\Gamma$ έχουμε $AB=3$, $A\Gamma=4$ και $\hat{A}=60^\circ$. Να υπολογίσετε τα εσωτερικά γινόμενα:
\begin{enumerate}
    \item $\vec{AB}\cdot\vec{A\Gamma}$
    \item $\vec{BA}\cdot\vec{A\Gamma}$
    \item $\vec{AB}\cdot\vec{B\Gamma}$
    \item $\vec{A\Gamma}\cdot(\vec{AB}+\vec{A\Gamma})$
\end{enumerate}

\askhsh
Δίνεται τετράγωνο $AB\Gamma\Delta$ πλευράς $a$. Να υπολογίσετε τα εσωτερικά γινόμενα:
\begin{enumerate}
    \item $\vec{AB}\cdot\vec{A\Delta}$
    \item $\vec{AB}\cdot\vec{A\Gamma}$
    \item $\vec{A\Gamma}\cdot\vec{B\Delta}$
    \item $\vec{AB}\cdot\vec{\Gamma\Delta}$
\end{enumerate}

\askhsh
Δίνονται τα μοναδιαία διανύσματα $\vec{i}, \vec{j}$ των αξόνων $x'x$ και $y'y$ αντίστοιχα. Να υπολογίσετε:
\begin{enumerate}
    \item $\vec{i}\cdot\vec{i}$
    \item $\vec{i}\cdot\vec{j}$
    \item $(\vec{i}+\vec{j})\cdot(\vec{i}-\vec{j})$
    \item $|2\vec{i}+3\vec{j}|^2$
\end{enumerate}

\paragraph{Εσωτερικό γινόμενο (συντεταγμένες)}

\askhsh
Δίνονται τα διανύσματα $\vec{a}=(3, -4)$ και $\vec{b}=(2, 1)$. Να βρείτε:
\begin{enumerate}
    \item Το εσωτερικό γινόμενο $\vec{a}\cdot\vec{b}$.
    \item Το μέτρο του $\vec{a}$.
    \item Το εσωτερικό γινόμενο $(2\vec{a})\cdot\vec{b}$.
    \item Το $\vec{a}\cdot(\vec{a}+\vec{b})$.
\end{enumerate}

\askhsh
Θεωρούμε τα σημεία $A(1, 2)$, $B(-1, 3)$ και $\Gamma(0, -2)$. Να υπολογίσετε:
\begin{enumerate}
    \item Τις συντεταγμένες των διανυσμάτων $\vec{AB}$ και $\vec{A\Gamma}$.
    \item Το εσωτερικό γινόμενο $\vec{AB}\cdot\vec{A\Gamma}$.
    \item Το εσωτερικό γινόμενο $\vec{BA}\cdot\vec{B\Gamma}$.
    \item Το μήκος της διαμέσου $\mu_a$ (υπόδειξη: βρείτε το μέσον $M$ και το $\vec{AM}^2$).
\end{enumerate}

\askhsh
Δίνονται τα διανύσματα $\vec{u}=(x, 2)$ και $\vec{v}=(3, -x)$.
\begin{enumerate}
    \item Να εκφράσετε το $\vec{u}\cdot\vec{v}$ ως συνάρτηση του $x$.
    \item Να βρείτε το $x$ ώστε $\vec{u}\cdot\vec{v}=1$.
    \item Να βρείτε το $x$ ώστε $\vec{u}\cdot\vec{v}=6$.
    \item Για $x=1$, να βρείτε το $(x\vec{u})\cdot\vec{v}$.
\end{enumerate}

\askhsh
Δίνονται τα διανύσματα $\vec{a}=(2, \lambda)$ και $\vec{b}=(\lambda-1, 3)$. Να βρείτε την τιμή του $\lambda \in \mathbb{R}$ ώστε:
\begin{enumerate}
    \item Το εσωτερικό γινόμενο να είναι 0.
    \item Το εσωτερικό γινόμενο να είναι 5.
    \item Τα διανύσματα να είναι κάθετα.
    \item Να ισχύει $\vec{a}\cdot\vec{b} = |\vec{a}|^2$ (για κάποια απλή τιμή).
\end{enumerate}

\paragraph{Ιδιότητες εσωτερικού γινομένου}

\askhsh
Αν για τα διανύσματα $\vec{a}, \vec{b}$ ισχύει $|\vec{a}|=3, |\vec{b}|=2$ και $(\widehat{\vec{a},\vec{b}})=\frac{2\pi}{3}$, να υπολογίσετε:
\begin{enumerate}
    \item Το $\vec{a}\cdot\vec{b}$.
    \item Το $(\vec{a}+\vec{b})^2$.
    \item Το $(2\vec{a}-\vec{b})^2$.
    \item Το $(3\vec{a}+2\vec{b})\cdot(\vec{a}-2\vec{b})$.
\end{enumerate}

\askhsh
Δίνονται $|\vec{u}|=2, |\vec{v}|=3$ και $\vec{u}\cdot\vec{v}=-1$. Να υπολογίσετε:
\begin{enumerate}
    \item Το $|\vec{u}+\vec{v}|^2$.
    \item Το $|\vec{u}-\vec{v}|^2$.
    \item Την παράσταση $\Pi = (\vec{u}+\vec{v})^2 - (\vec{u}-\vec{v})^2$.
    \item Το γινόμενο $(\vec{u}+2\vec{v})\cdot(2\vec{u}-\vec{v})$.
\end{enumerate}

\askhsh
Αν $|\vec{a}|=|\vec{b}|=1$ και $\vec{a}\cdot\vec{b}=0$, να αποδείξετε ότι:
\begin{enumerate}
    \item $|\vec{a}+\vec{b}| = \sqrt{2}$.
    \item $|\vec{a}-\vec{b}| = \sqrt{2}$.
    \item $(\vec{a}+\vec{b})\cdot(\vec{a}-\vec{b}) = 0$.
    \item Η γωνία των $\vec{a}+\vec{b}$ και $\vec{a}$ είναι $45^\circ$.
\end{enumerate}

\askhsh
Αν για τα διανύσματα $\vec{x}, \vec{y}$ ισχύει $2\vec{x}+3\vec{y}=\vec{0}$ και $|\vec{x}|=3$, να υπολογίσετε:
\begin{enumerate}
    \item Το μέτρο $|\vec{y}|$.
    \item Το εσωτερικό γινόμενο $\vec{x}\cdot\vec{y}$.
    \item Το γινόμενο $(\vec{x}+\vec{y})\cdot\vec{y}$.
    \item Τη γωνία $(\widehat{\vec{x},\vec{y}})$.
\end{enumerate}

\paragraph{Κάθετα διανύσματα}

\askhsh 
Να εξετάσετε αν τα παρακάτω ζεύγη διανυσμάτων είναι κάθετα:
\begin{enumerate}
    \item $\vec{a}=(1, -2)$ και $\vec{b}=(4, 2)$.
    \item $\vec{a}=(2, 3)$ και $\vec{b}=(-3, 2)$.
    \item $\vec{a}=(\sqrt{3}, 1)$ και $\vec{b}=(1, \sqrt{3})$.
    \item $\vec{a}=(x, y)$ και $\vec{b}=(-y, x)$.
\end{enumerate}

\askhsh
Δίνεται τρίγωνο με κορυφές $A(2, 1)$, $B(5, 1)$ και $\Gamma(2, 4)$.
\begin{enumerate}
    \item Να βρείτε τα διανύσματα $\vec{AB}, \vec{A\Gamma}$.
    \item Να δείξετε ότι $\vec{AB} \perp \vec{A\Gamma}$.
    \item Να βρείτε το $\vec{B\Gamma}$.
    \item Να επιβεβαιώσετε το Πυθαγόρειο Θεώρημα για τα μέτρα των πλευρών.
\end{enumerate}

\askhsh
Έστω $\vec{u}, \vec{v}$ διανύσματα με $|\vec{u}|=|\vec{v}|=1$. Να δείξετε ότι τα διανύσματα $\vec{a} = \vec{u} + \vec{v}$ και $\vec{b} = \vec{u} - \vec{v}$ είναι κάθετα, υπολογίζοντας:
\begin{enumerate}
    \item Το εσωτερικό γινόμενο $\vec{a}\cdot\vec{b}$.
    \item Αναπτύσσοντας το $(\vec{u}+\vec{v})(\vec{u}-\vec{v})$.
    \item Το $|\vec{u}|^2 - |\vec{v}|^2$.
    \item Τι συμπέρασμα βγάζετε για τις διαγωνίους ρόμβου;
\end{enumerate}

\askhsh
Να βρείτε τον αριθμό $x \in \mathbb{R}$ ώστε τα διανύσματα να είναι κάθετα:
\begin{enumerate}
    \item $\vec{a}=(x, 2)$ και $\vec{b}=(1, -3)$.
    \item $\vec{u}=(x-1, 2)$ και $\vec{v}=(x, -3)$.
    \item $\vec{a}=(\sin x, \cos x)$ και $\vec{b}=(\cos x, -\sin x)$.
    \item $\vec{w}=(2, x)$ και $\vec{z}=(x, -8)$ (όπου $x>0$).
\end{enumerate}

\paragraph{Υπολογισμός μέτρου γραμμικού συνδυασμού διανυσμάτων}

\askhsh
Δίνονται $|\vec{a}|=2, |\vec{b}|=1, (\widehat{\vec{a},\vec{b}})=60^\circ$. Να υπολογίσετε τα μέτρα:
\begin{enumerate}
    \item $|\vec{a}+\vec{b}|$.
    \item $|\vec{a}-\vec{b}|$.
    \item $|2\vec{a}+\vec{b}|$.
    \item $|\vec{a}-2\vec{b}|$.
\end{enumerate}

\askhsh
Αν $|\vec{u}|=2, |\vec{v}|=2\sqrt{3}, \vec{u}\cdot\vec{v}=-6$, να υπολογίσετε:
\begin{enumerate}
    \item Το $\vec{u}^2 + \vec{v}^2$.
    \item Το $(\vec{u}+\vec{v})^2$.
    \item Το μέτρο $|\vec{u}+\vec{v}|$.
    \item Το μέτρο $|\vec{u}-\vec{v}|$.
\end{enumerate}

\askhsh
Δίνονται τα διανύσματα $\vec{a}, \vec{b}$ με $|\vec{a}|=3, |\vec{b}|=5$ και $|\vec{a}+\vec{b}|=7$. Να βρείτε:
\begin{enumerate}
    \item Το τετράγωνο $|\vec{a}+\vec{b}|^2$.
    \item Το εσωτερικό γινόμενο $\vec{a}\cdot\vec{b}$.
    \item Το $|\vec{a}-\vec{b}|^2$.
    \item Το μέτρο $|\vec{a}-\vec{b}|$.
\end{enumerate}

\askhsh
Έστω ισόπλευρο τρίγωνο $AB\Gamma$ πλευράς 4. Να υπολογίσετε τα μέτρα των διανυσμάτων:
\begin{enumerate}
    \item $|\vec{AB}+\vec{A\Gamma}|$.
    \item $|\vec{AB}-\vec{A\Gamma}|$.
    \item $|2\vec{AB}-\vec{A\Gamma}|$.
    \item $|\vec{AB}-2\vec{B\Gamma}|$.
\end{enumerate}

\paragraph{Συνημίτονο γωνίας διανυσμάτων}

\askhsh
Δίνονται τα διανύσματα $\vec{a}=(1, 2)$ και $\vec{b}=(3, 1)$. Να υπολογίσετε:
\begin{enumerate}
    \item Το $\vec{a}\cdot\vec{b}$.
    \item Τό μέτρο $|\vec{a}|$.
    \item Το μέτρο $|\vec{b}|$.
    \item Το συνημίτονο της γωνίας $(\widehat{\vec{a},\vec{b}})$.
\end{enumerate}

\askhsh
Αν $|\vec{u}|=4, |\vec{v}|=2$ και $\vec{u}\cdot\vec{v}=4\sqrt{2}$, να βρείτε:
\begin{enumerate}
    \item Το $\cos(\widehat{\vec{u},\vec{v}})$.
    \item Τη γωνία $(\widehat{\vec{u},\vec{v}})$.
    \item Τη γωνία $(\widehat{\vec{u}, -\vec{v}})$.
    \item Τη γωνία $(\widehat{-\vec{u}, -\vec{v}})$.
\end{enumerate}

\askhsh
Δίνεται τρίγωνο με κορυφές $A(2,3), B(-1, 0), \Gamma(3, -2)$. Να υπολογίσετε:
\begin{enumerate}
    \item Τα διανύσματα $\vec{AB}, \vec{A\Gamma}$.
    \item Το γινόμενο $\vec{AB}\cdot\vec{A\Gamma}$ και τα μέτρα τους.
    \item Το συνημίτονο της γωνίας $\hat{A}$.
    \item Το είδος της γωνίας $\hat{A}$ (οξεία, αμβλεία, ορθή).
\end{enumerate}

\askhsh
Δίνονται τα διανύσματα $\vec{a}, \vec{b}$ με $2|\vec{a}|=|\vec{b}|$ και $|\vec{a}+\vec{b}| = \sqrt{7}|\vec{a}|$. Να βρείτε:
\begin{enumerate}
    \item Την σχέση $|\vec{b}|^2$ και $|\vec{a}|^2$.
    \item Την αντικατάσταση στη σχέση του μέτρου αθροίσματος.
    \item Το συνημίτονο της γωνίας τους.
    \item Το μέτρο της γωνίας τους.
\end{enumerate}

\paragraph{Αποδεικτικές ασκήσεις}

\askhsh
Αν για τα διανύσματα $\vec{u}, \vec{v}$ ισχύει $|\vec{u}|=|\vec{v}|$, να αποδείξετε ότι:
\begin{enumerate}
    \item $(\vec{u}+\vec{v})\cdot(\vec{u}-\vec{v}) = 0$.
    \item Τα διανύσματα $\vec{u}+\vec{v}$ και $\vec{u}-\vec{v}$ είναι κάθετα.
    \item Το άθροισμα $\vec{u}+\vec{v}$ διχοτομεί τη γωνία των $\vec{u}, \vec{v}$.
    \item Αν επιπλέον $|\vec{u}+\vec{v}|=|\vec{u}|$, τότε η γωνία είναι $120^\circ$.
\end{enumerate}

\askhsh
Να αποδείξετε τις παρακάτω ταυτότητες για τυχαία διανύσματα:
\begin{enumerate}
    \item $|\vec{a}+\vec{b}|^2 + |\vec{a}-\vec{b}|^2 = 2(|\vec{a}|^2 + |\vec{b}|^2)$.
    \item $|\vec{u}+\vec{v}|^2 - |\vec{u}-\vec{v}|^2 = 4\vec{u}\cdot\vec{v}$.
    \item Αν $\vec{a} \perp \vec{b}$, τότε $|\vec{a}+\vec{b}|^2 = |\vec{a}|^2 + |\vec{b}|^2$.
    \item $ (\alpha \vec{u}) \cdot (\beta \vec{v}) = \alpha \beta (\vec{u} \cdot \vec{v}) $.
\end{enumerate}

\askhsh
Αν $\vec{a} + \vec{b} + \vec{c} = \vec{0}$ και $|\vec{a}|=3, |\vec{b}|=4, |\vec{c}|=5$:
\begin{enumerate}
    \item Να αποδείξετε ότι $\vec{a}+\vec{b} = -\vec{c}$.
    \item Να βρείτε το $\vec{a}\cdot\vec{b}$.
    \item Να δείξετε ότι $\vec{a} \perp \vec{b}$.
    \item Να ερμηνεύσετε το αποτέλεσμα γεωμετρικά.
\end{enumerate}

\askhsh
Δίνονται τα μη μηδενικά διανύσματα $\vec{a}, \vec{b}$. Να αποδείξετε ότι:
\begin{enumerate}
    \item Αν $|\vec{a}+\vec{b}| = |\vec{a}| + |\vec{b}|$, τότε $\vec{a} \uparrow\uparrow \vec{b}$.
    \item Αν $|\vec{a}+\vec{b}| = ||\vec{a}| - |\vec{b}||$, τότε $\vec{a} \uparrow\downarrow \vec{b}$.
    \item Αν $|\vec{a}+\vec{b}| = |\vec{a}-\vec{b}|$, τότε $\vec{a} \perp \vec{b}$.
    \item Αν $\vec{a}\cdot\vec{b} = |\vec{a}||\vec{b}|$, τότε $\vec{a} \uparrow\uparrow \vec{b}$.
\end{enumerate}

\paragraph{Εύρεση παραμέτρων}

\askhsh
Δίνονται τα διανύσματα $\vec{a}=(\lambda, 2)$ και $\vec{b}=(-1, \lambda)$. Να βρείτε το $\lambda$ ώστε:
\begin{enumerate}
    \item Τα διανύσματα να είναι κάθετα.
    \item Τα διανύσματα να είναι παράλληλα.
    \item Να ισχύει $\vec{a}\cdot\vec{b} = 3$.
    \item Το μέτρο του $\vec{a}$ να είναι ίσο με $\sqrt{5}$.
\end{enumerate}

\askhsh
Έστω $\vec{u}=(1, \lambda)$ και $\vec{v}=(\lambda-2, 3)$.
\begin{enumerate}
    \item Για ποια τιμή του $\lambda$ είναι $\vec{u} \perp \vec{v}$;
    \item Για την τιμή που βρήκατε, να υπολογίσετε το $|\vec{u}+\vec{v}|$.
    \item Να βρείτε το $\lambda$ ώστε $|\vec{u}|=\sqrt{5}$.
    \item Αν $\lambda=2$, να βρείτε τη γωνία των διανυσμάτων.
\end{enumerate}

\askhsh
Δίνονται $\vec{a}, \vec{b}$ με $|\vec{a}|=1$, $|\vec{b}|=2$. Ζητείται ο $k \in \mathbb{R}$ ώστε το διάνυσμα $\vec{w} = \vec{a} + k\vec{b}$ να είναι κάθετο στο $\vec{a}$.
\begin{enumerate}
    \item Να γράψετε τη συνθήκη καθετότητας με εσωτερικό γινόμενο.
    \item Αν επιπλέον $\phi=60^\circ$, να βρείτε το $\vec{a}\cdot\vec{b}$.
    \item Να λύσετε την εξίσωση ως προς $k$.
    \item Να βρείτε το μέτρο του $\vec{w}$ για αυτό το $k$.
\end{enumerate}

\askhsh
Θεωρούμε $\vec{p} = \kappa \vec{a} + \vec{b}$ και $\vec{q} = 2\vec{a} - \vec{b}$, όπου $\vec{a}, \vec{b}$ κάθετα μοναδιαία διανύσματα.
\begin{enumerate}
    \item Να βρείτε το $\vec{p}\cdot\vec{q}$ συναρτήσει του $\kappa$.
    \item Να βρείτε το $\kappa$ ώστε $\vec{p} \perp \vec{q}$.
    \item Να βρείτε το $\kappa$ ώστε $\vec{p} // \vec{q}$.
    \item Να βρείτε το $\kappa$ ώστε $|\vec{p}| = \sqrt{5}$.
\end{enumerate}

\end{document}
